\section{Primitive-grid determinants, local defects and finite descent}\label{jets:sec:det}
The polynomial character normal form above remains valid over arbitrary
complex coefficient algebras. Here identify the prime weights $A_p$ with
$t_p$, write $\mathcal D_q$ for that quotient, and put
$v_{m,\chi}=Y_{m,\chi}$, $b_\chi=Y_{f,\chi}$ for
$f=\operatorname{cond}(\chi)$. The coefficient $E_{m,\chi}$ is
$A_{m,\chi}$ of Theorem~\ref{dist:thm:dist-normal-form}. Thus
\begin{equation}\label{jets:eq:conductor-normal}
 [v_{m,\chi}]=E_{m,\chi}b_\chi.
\end{equation}
No Euler-factor division is used.
Let $U_q=\{u\bmod q:(u,q)=1\}$, and use $r=0$ for the Hurwitz endpoint
$x_r=1$. Distinguish the rank of the full quotient from the image of
its primitive-grid coordinate map.

Conductor coordinates are polynomial and globally valid, but it is often
more convenient to use the $\varphi(q)$ primitive-grid points
$e_{u/q}$, $u\in U_q$. Fourier-transforming those points gives
$v_{q,\chi}$. By Theorem~\ref{dist:thm:dist-normal-form}, the map from their formal
span into $\D_q$ is diagonal in character coordinates:
\begin{equation}\label{jets:eq:primitive-diagonal}
 v_{q,\chi}\longmapsto E_{q,\chi}b_\chi.
\end{equation}
Its determinant in these specified Fourier and conductor bases is
$\Delta_q=\prod_{\chi\bmod q}E_{q,\chi}$.

For $p^e\Vert q$, write
\[
 M_p=q/p^e,\qquad h_p=\ord_{M_p}(p),
\]
with $h_p=1$ when $M_p=1$.

\begin{theorem}[Determinant factorization]\label{jets:thm:determinant}
One has the polynomial identity
\begin{equation}\label{jets:eq:determinant}
 \boxed{\displaystyle
 \Delta_q(\mathbf A)=
 \prod_{p^e\Vert q}
 A_p^{\varphi(M_p)(p^{e-1}-1)}
 \bigl(A_p^{h_p}-1\bigr)^{\varphi(M_p)/h_p}.}
\end{equation}
Thus primitive-grid coordinates form a basis precisely when all the
factors in \eqref{jets:eq:determinant} are invertible in a field
specialization. Their failure does not change the rank of $\D_q$.
\end{theorem}
\begin{proof}
Fix $p^e\Vert q$ and inspect the $p$ factor of
$\prod_\chi E_{q,\chi}$. Characters whose conductor is not divisible by
$p$ are exactly the characters of $U_{M_p}$. Their additional factors
are $A_p-\chi(p)$. In the character group of $U_{M_p}$, the values
$\chi(p)$ run through every $h_p$th root of unity with multiplicity
$\varphi(M_p)/h_p$. This follows by duality from the fact that $p$ has
order $h_p$ in $U_{M_p}$. Their product is therefore
$(A_p^{h_p}-1)^{\varphi(M_p)/h_p}$.

It remains to count the pure powers of $A_p$. If $a=v_p(f)$, the
exponent contributed by one character is $e-\max(1,a)$. Writing this
as a sum of indicators gives
\[
 \sum_{\chi\bmod q}\bigl(e-\max(1,v_p(f_\chi))\bigr)
 =\sum_{j=1}^{e-1}\#\{\chi:v_p(f_\chi)\le j\}.
\]
The count in the last sum is $\varphi(p^j)\varphi(M_p)$, since these
are exactly the characters induced from level $p^jM_p$. Hence the
exponent is
$\varphi(M_p)\sum_{j=1}^{e-1}\varphi(p^j)
=\varphi(M_p)(p^{e-1}-1)$. Multiply over the primes.
\end{proof}

\begin{example}[Level twelve]
Let $\chi_3$ and $\chi_4$ be the nonprincipal quadratic characters of
conductors three and four, and let $\chi_{12}=\chi_3\chi_4$.
Their conductor generators are
\begin{align*}
 b_3&=e_{1/3}-e_{2/3},&
 b_4&=e_{1/4}-e_{3/4},\\
 b_{12}&=e_{1/12}-e_{5/12}-e_{7/12}+e_{11/12}.
\end{align*}
For $r\in U_{12}$, the normal form reads
\begin{align}\label{jets:eq:q12-normal}
 [e_{r/12}]=\frac14\bigl[&A_2(A_2-1)(A_3-1)e_0
 +\chi_3(r)A_2(A_2+1)b_3\notag\\
 &+\chi_4(r)(A_3+1)b_4+\chi_{12}(r)b_{12}\bigr].
\end{align}
Consequently
\[
 \Delta_{12}=A_2^2(A_2^2-1)(A_3^2-1).
\]
At $A_2=A_3=1$, the primitive-grid span misses the principal conductor
coordinate, even though the full quotient still has rank four.
\end{example}

\subsection{Local jet defects of the primitive-grid map}
Specialize to $A_p(s)=p^s$, and let $s=s_0+t$. All exponential factors
are nonzero. A factor $p^s-\chi(p)$ vanishes precisely when
$p^{s_0}=\chi(p)$, and any such zero is simple because its derivative
is $p^{s_0}\log p\ne0$.

\begin{theorem}[Local Smith factors and primitive-jet ranks]\label{jets:thm:smith}
Define
\begin{equation}\label{jets:eq:rho-general}
 \rho_\chi(s_0)=\#\{p\mid q:p\nmid f_\chi,\ p^{s_0}=\chi(p)\}.
\end{equation}
Over $\C[[t]]$, the diagonal factors of the primitive-grid map are
units times $t^{\rho_\chi(s_0)}$. Consequently its rank on order-$N$
truncated jets, as a complex linear map, is
\begin{equation}\label{jets:eq:primitive-jet-rank}
 \sum_{\chi\bmod q}\max\{N+1-\rho_\chi(s_0),0\}.
\end{equation}
At $s_0=0$, this is determined by
$\rho_\chi(0)=\#\{p\mid q:p\nmid f_\chi,\ \chi(p)=1\}$.
\end{theorem}
\begin{proof}
Factor $E_{q,\chi}(s_0+t)$ into its simple vanishing factors and a unit.
The primitive Fourier matrix is constant and invertible, so these are
also the local invariant factors, up to permutation and units.
Multiplication by $t^r$ on $\C[t]/(t^{N+1})$ has rank
$\max(N+1-r,0)$. Summing over the diagonal coordinates proves the
formula. This is the rank of a \emph{coordinate inclusion map}, not
the rank of the distribution relation matrix.
\end{proof}

\begin{corollary}[Multiple-prime resonances are central]\label{jets:cor:central}
If $s_0\ne0$, at most one prime can satisfy
$p^{s_0h_p}=1$. At a nonzero resonant point for that prime $p$, the
primitive-grid rank defect is $\varphi(M_p)/h_p$, and every affected
local factor has order one. The order-$N$ inclusion rank is then
$(N+1)\varphi(q)-\varphi(M_p)/h_p$ for all $N\ge0$.
\end{corollary}
\begin{proof}
A resonance forces $s_0$ to be purely imaginary. Resonances at distinct
primes would make $\log p/\log\ell$ rational, because each phase is a
rational multiple of $2\pi$. Unique prime factorization rules this out
unless $s_0=0$. For the one resonant prime, the multiplicity count in
Theorem~\ref{jets:thm:determinant} gives the number of affected characters.
Apply Theorem~\ref{jets:thm:smith}.
\end{proof}

\begin{example}[Two exact jet-rank profiles]
At level $12$ and $s_0=0$, the principal character has $\rho=2$ and the
other three characters have $\rho=0$. For $N=0,1,2$, the primitive-jet
ranks are $3,6,10$, while the full quotient dimensions are $4,8,12$.
At level $21$, the principal character has $\rho=2$, the character
induced from conductor three has $\rho=1$, and the other ten
characters have $\rho=0$. The inclusion ranks are $10,21,33$, compared
with full quotient dimensions $12,24,36$. Indeed,
\[
 \Delta_{21}=(A_3^6-1)(A_7-1)^2.
\]
Exact character-phase certificates for these examples are supplied.
\end{example}

\section{Explicit primitive-grid reduction coefficients}\label{jets:sec:primitive}

Away from the exceptional locus, the conductor normal form can be
inverted to express all grid points through the primitive ones. For
Hurwitz weights there is a particularly concrete answer that avoids
character-field arithmetic.

Let $\mathcal S_q$ be the multiplicative semigroup of positive integers
whose prime factors divide $q$, including $1$. For $\Real s>0$, define
\begin{equation}\label{jets:eq:C-smooth}
 C_{r,u}(s)=\sum_{\substack{d\in\mathcal S_q\\du\equiv r\; (q)}}d^{-s},
 \qquad 0\le r<q,\quad u\in U_q.
\end{equation}
The series is absolutely convergent, since
$\sum_{d\in\mathcal S_q}d^{-\sigma}
=\prod_{p\mid q}(1-p^{-\sigma})^{-1}$ for $\sigma>0$.

\begin{theorem}[Primitive-grid extension]\label{jets:thm:primitive}
For $\Real s>0$, the $q\times\varphi(q)$ matrix $C(s)$ satisfies every
prime and divisor distribution relation with weights $p^s$, and its
primitive rows form the identity matrix. It is the unique extension
matrix from primitive-grid data. In particular,
\begin{equation}\label{jets:eq:Hurwitz-C}
 \zeta(s,r/q)=\sum_{u\in U_q}C_{r,u}(s)\zeta(s,u/q)
\end{equation}
with the endpoint convention for $r=0$, initially for $\Real s>1$ and
then meromorphically wherever the coefficients are defined.
\end{theorem}
\begin{proof}
For a primitive row $r$, the congruence $du\equiv r\pmod q$ forces
$d$ to be prime to $q$, hence $d=1$. This gives the identity submatrix.
For a prime distribution row indexed by $a\equiv0\pmod p$,
\[
 \sum_{pr\equiv a\;(q)}C_{r,u}(s)
 =\sum_{\substack{d\in\mathcal S_q\\pdu\equiv a\;(q)}}d^{-s}
 =p^s\sum_{\substack{d'\in\mathcal S_q\\d'u\equiv a\;(q)}}(d')^{-s}.
\]
In the last step $d'=pd$; conversely every $d'$ in the last sum is
divisible by $p$, because $u$ is a unit and $a$ is divisible by $p$.
The last expression is $p^sC_{a,u}(s)$. Prime sufficiency gives all
divisor rows. Theorem~\ref{dist:thm:dist-normal-form} and the identity submatrix
show that these $\varphi(q)$ independent columns exhaust the solution
space. Hurwitz zeta satisfies the same distribution equations by
regrouping its absolutely convergent defining series, proving
\eqref{jets:eq:Hurwitz-C}.
\end{proof}

\subsection{A finite formula, including negative spectral parameters}
For $0\le r<q$, let
\[
 g=\gcd(r,q),\quad m=q/g,\quad a=r/g,
 \quad P_r=\{p\mid q:p\nmid m\}.
\]
When $r=0$, this means $g=q$, $m=1$, $a=0$.
For $p\in P_r$, put $h_{p,m}=\ord_m(p)$, with order one modulo one.

\begin{theorem}[Finite residue formula]\label{jets:thm:finite-C}
The meromorphic continuation of \eqref{jets:eq:C-smooth} is
\begin{equation}\label{jets:eq:C-finite}
 C_{r,u}(s)=g^{-s}
 \frac{\displaystyle
 \sum_{\substack{0\le b_p<h_{p,m}\;(p\in P_r)\\
 u\prod_{p\in P_r}p^{b_p}\equiv a\;(m)}}
 \prod_{p\in P_r}p^{-s b_p}}
 {\displaystyle\prod_{p\in P_r}(1-p^{-s h_{p,m}})}.
\end{equation}
Empty products have their usual value one. In particular,
\begin{equation}\label{jets:eq:C-endpoint}
 C_{0,u}(s)=q^{-s}\prod_{p\mid q}(1-p^{-s})^{-1},
\end{equation}
independently of $u$.
\end{theorem}
\begin{proof}
The congruence in \eqref{jets:eq:C-smooth} forces the valuation of $d$ at
an unsaturated prime, one dividing $m$, to be exactly $v_p(g)$.
At a saturated prime, one not dividing $m$, it is at least $v_p(g)$.
Thus $d=g\prod_{p\in P_r}p^{c_p}$ with $c_p\ge0$, and the remaining
congruence is $u\prod p^{c_p}\equiv a\pmod m$.
Write $c_p=b_p+h_{p,m}j_p$ and sum the independent geometric series in
$j_p$. This proves the formula in $\Real s>0$ and hence its meromorphic
continuation. The endpoint has trivial congruence modulo one.
\end{proof}

Formula \eqref{jets:eq:C-finite}, \emph{not} the divergent smooth-number
series, is the definition used at negative real $s$. It has no pole at
any nonzero real $s$, and the continued distribution identities remain
valid there. At a zero of a displayed denominator, one should use the
polynomial conductor normal form rather than cancel poles numerically.
Some individual coefficients may have removable singularities, but
Theorem~\ref{jets:thm:determinant} gives the criterion for the entire
primitive-grid coordinate system.

\begin{corollary}[Exact support and signs]\label{jets:cor:support}
Let $H_r$ be the subgroup of $U_m$ generated by the classes of the
primes in $P_r$. For every nonzero real $s$,
\[
 C_{r,u}(s)\ne0\quad\Longleftrightarrow\quad au^{-1}\in H_r.
\]
The number of supported primitive coordinates is
\begin{equation}\label{jets:eq:support-count}
 \frac{\varphi(q)}{\varphi(m)}|H_r|.
\end{equation}
All supported coefficients are positive for $s>0$. For $s<0$, they
all have sign $(-1)^{|P_r|}$.
\end{corollary}
\begin{proof}
The numerator of \eqref{jets:eq:C-finite} is a sum of positive real terms
and is nonempty exactly under the subgroup condition. The reduction
map $U_q\to U_m$ has fibers of cardinality
$\varphi(q)/\varphi(m)$, giving \eqref{jets:eq:support-count}.
For $s>0$ every denominator factor is positive. For $s<0$ each of the
$|P_r|$ denominator factors is negative, while the numerator and
$g^{-s}$ are positive.
\end{proof}

\begin{corollary}[Complete monotonicity and normalization]\label{jets:cor:monotone}
For $s>0$ and $j\ge0$, put
\begin{equation}\label{jets:eq:W-definition}
 W_{r,u}^{[j]}(s)=(-1)^j\frac{d^j}{ds^j}C_{r,u}(s)
 =\sum_{\substack{d\in\mathcal S_q\\du\equiv r\;(q)}}
 \frac{(\log d)^j}{d^s}.
\end{equation}
These quantities are nonnegative. Moreover,
$\sum_u C_{r,u}(1)=1$ for every $r$, and
$0\le C_{r,u}(k)\le1$ for integer $k\ge1$.
\end{corollary}
\begin{proof}
The logarithmically weighted semigroup series converges locally
uniformly in $\Real s>0$, so termwise differentiation is valid.
The constant grid vector satisfies the distribution equations at
$s=1$, and its primitive entries are all one. Uniqueness of extension
proves the row-sum identity. Monotonicity gives
$C_{r,u}(k)\le C_{r,u}(1)\le1$.
\end{proof}

\begin{example}[Two reductions at denominator six]\label{jets:ex:q6}
The nontrivial coefficients for $r=2$ are
\[
 C_{2,1}(s)=\frac{2^s}{2^{2s}-1},\qquad
 C_{2,5}(s)=\frac1{2^{2s}-1}.
\]
Thus, meromorphically,
\begin{align}
 \zeta(s,1/3)&=
 \frac{2^s\zeta(s,1/6)+\zeta(s,5/6)}{2^{2s}-1},\label{jets:eq:q6-one}\\
 \zeta(s,1/2)&=
 \frac{\zeta(s,1/6)+\zeta(s,5/6)}{3^s-1}.
 \label{jets:eq:q6-two}
\end{align}
At $s=1$, the coefficients in \eqref{jets:eq:q6-one} are $2/3$ and $1/3$.
At $s=-1$ they are $-2/3$ and $-4/3$, illustrating the sign theorem
without using a divergent series. The poles at $s=0$ reflect a
singular primitive basis, not a failure of the underlying distribution
relations.
\end{example}

\subsection{Denominator certificates and construction cost}
Let $M_p$ also denote, in this paragraph only, the pushforward matrix
on grid vectors defined by
\[
 (M_pf)_r=\sum_{pu\equiv r\;(q)}f_u.
\]
To avoid confusion with the integer complement in
Section~\ref{jets:sec:det}, write that complement as $q_p=q/p^e$ here and
put $h_p=\ord_{q_p}(p)$. Multiplication of residues gives
$M_p^{e+h_p}=M_p^e$. For $z=p^{-s}$,
\begin{equation}\label{jets:eq:operator-inverse}
 (I-zM_p)^{-1}
 =\sum_{j=0}^{e-1}z^jM_p^j
 +\frac{z^e}{1-z^{h_p}}\sum_{j=0}^{h_p-1}z^jM_p^{e+j}.
\end{equation}
This follows first from a geometric series for $|z|<1$, then as a
rational matrix identity. The commuting product of these inverses,
applied to vectors supported on $U_q$, is the extension $C(s)$.

\begin{proposition}[A uniform rational denominator]\label{jets:prop:denominator}
For integer $k\ge1$, every reduced denominator of $C_{r,u}(k)$ divides
\begin{equation}\label{jets:eq:denominator-bound}
 D_q(k)=\prod_{p^e\Vert q}
 p^{k(e-1)}\bigl(p^{kh_p}-1\bigr).
\end{equation}
Its bit length is $O(kq\log q)$. The coefficients of the raw derivative
$C_{r,u}^{(j)}(k)$, viewed as a polynomial in the symbols $\log p$,
have denominators dividing $D_q(k)^{j+1}$.
\end{proposition}
\begin{proof}
Every scalar coefficient on the right of
\eqref{jets:eq:operator-inverse} has denominator dividing
$p^{k(e-1)}(p^{kh_p}-1)$. Multiplication over primes proves the first
claim. Since $h_p\le q$ and $\sum_{p\mid q}\log p\le\log q$, its
logarithm is at most a constant times $kq\log q$.
The matrix product can be written as a ratio of integer polynomials in
$p^s$, with denominator obtained from \eqref{jets:eq:denominator-bound} by
replacing $k$ with $s$. Repeated differentiation introduces polynomial
factors in $\log p$ and raises the denominator power by at most one
each time. For normalized Taylor coefficients, an additional factor
$j!$ may be needed.
\end{proof}

The implementation computes \eqref{jets:eq:C-finite} by convolution on
residue classes modulo $m$, separately for each $m\mid q$. At a fixed
$m$, each prime update takes at most $\varphi(m)^2$ rational arithmetic
operations. A complete matrix can therefore be constructed in
\[
 O\!\left(\omega(q)\sum_{m\mid q}\varphi(m)^2+q\varphi(q)\right)
 =O\bigl(\omega(q)q^2\bigr)
\]
rational arithmetic operations, using $\sum_{m\mid q}\varphi(m)=q$.
This is an arithmetic-operation bound, not a constant-cost bit model
or a claim of optimality. The denominator estimate supplies a
conservative size control at positive integer orders.

\section{All-orders Hurwitz--Stieltjes identities}\label{jets:sec:stieltjes}

Write the Laurent expansion as
\begin{equation}\label{jets:eq:stieltjes-laurent}
 \zeta(1+t,x)=\frac1t+
 \sum_{n\ge0}\frac{(-1)^n}{n!}\gamma_n(x)t^n,
 \qquad x>0.
\end{equation}
Here $\gamma_n^{(k)}$ always means $k$ derivatives in $x$; a derivative
on $C(s)$ is in the spectral variable. The coefficient identities in
this section are analytic theorems, distinct from formal rank claims.

\begin{theorem}[Derived Stieltjes reduction]\label{jets:thm:derived-stieltjes}
For $k\ge1$, $n\ge0$, and $x_r=r/q$ with the endpoint convention,
\begin{equation}\label{jets:eq:derived-stieltjes}
 \boxed{\displaystyle
 \gamma_n^{(k)}(x_r)=
 \sum_{u\in U_q}\sum_{j=0}^n\binom nj
 W_{r,u}^{[j]}(k+1)\,
 \gamma_{n-j}^{(k)}(u/q).}
\end{equation}
All the coefficient functions are given by the finite formula
\eqref{jets:eq:C-finite} and its derivatives.
\end{theorem}
\begin{proof}
Differentiating Hurwitz zeta in its parameter yields
\[
 \partial_x^k\zeta(1+t,x)=(-1)^k(1+t)_k\zeta(k+1+t,x).
\]
The pole in \eqref{jets:eq:stieltjes-laurent} is independent of $x$, so for
$k\ge1$ the left side is
$\sum_n(-1)^n\gamma_n^{(k)}(x)t^n/n!$.
Apply \eqref{jets:eq:Hurwitz-C} at spectral parameter $k+1+t$ and multiply
by the common scalar $(-1)^k(1+t)_k$. Taylor multiplication gives
\eqref{jets:eq:derived-stieltjes}, because
$(-1)^j C_{r,u}^{(j)}=W_{r,u}^{[j]}$.
\end{proof}

This proof also explains why the harmonic-number master identity in
Chapter 8 is compatible with the same reduction matrix at every
Stieltjes order: the rising factorial is common to every rational
argument. It does not create new independent distribution coordinates.

\begin{theorem}[Underived Stieltjes reduction, with the pole term]\label{jets:thm:underived}
For $n\ge0$,
\begin{equation}\label{jets:eq:underived-stieltjes}
 \boxed{\displaystyle
 \gamma_n(x_r)=
 \sum_{u\in U_q}\sum_{j=0}^n\binom nj
 W_{r,u}^{[j]}(1)\gamma_{n-j}(u/q)
 -\frac1{n+1}\sum_{u\in U_q}W_{r,u}^{[n+1]}(1).}
\end{equation}
\end{theorem}
\begin{proof}
Apply \eqref{jets:eq:Hurwitz-C} at $s=1+t$. The coefficients are analytic
there. Multiplying their Taylor expansions by
\eqref{jets:eq:stieltjes-laurent}, the coefficient of $t^n$ receives the
usual product terms and also
$\sum_u C_{r,u}^{(n+1)}(1)/(n+1)!$ from the pole $1/t$.
Multiplication by $(-1)^n n!$ converts that additional term to the
negative last term of \eqref{jets:eq:underived-stieltjes}. The coefficient
of $1/t$ agrees because $\sum_u C_{r,u}(1)=1$.
\end{proof}

Omitting the last term would already give an incorrect formula for
$\gamma_0$. The point is not a rank correction: a known scalar right
side does not change a homogeneous rank. It is a correction to the
actual analytic value identity.

\begin{example}[An explicit first-order identity]
From \eqref{jets:eq:q6-one}, direct differentiation at $s=1$ gives
\[
 W_{2,1}^{[1]}(1)=\frac{10}{9}\log2,\qquad
 W_{2,5}^{[1]}(1)=\frac89\log2.
\]
Consequently
\begin{equation}\label{jets:eq:gamma0-six}
 \gamma_0(1/3)=\frac23\gamma_0(1/6)+\frac13\gamma_0(5/6)-2\log2.
\end{equation}
At $s=2$, the analogous coefficients give
\begin{align}\label{jets:eq:gamma1prime-six}
 \gamma_1'(1/3)={}&\frac4{15}\gamma_1'(1/6)
 +\frac1{15}\gamma_1'(5/6)\notag\\
 &+\frac{68\log2}{225}\gamma_0'(1/6)
 +\frac{32\log2}{225}\gamma_0'(5/6).
\end{align}
Since $\gamma_0'=-\psi'$, the signs can also be checked against the
parameter-differentiated digamma multiplication identity.
The verification program checks \eqref{jets:eq:underived-stieltjes} through
$n=2$ and \eqref{jets:eq:gamma1prime-six} by numerical parameter
differentiation of Stieltjes functions.
\end{example}

\section{Finite conductor descent for all Stieltjes layers}

For a primitive character $\chi$ of conductor $f\mid d$, define
\begin{equation}
G_{d,\chi}^{(n,k)}=
\sum_{a\in\U_d}\chi(a)\gamma_n^{(k)}(a/d).
\end{equation}
Set
\begin{equation}\label{conductor:eq:Rdf}
R(d,f)=\prod_{\substack{p\mid d\\p\nmid f}}p,
\qquad m_r=\frac{d}{fr}\quad(r\mid R(d,f)).
\end{equation}
Each $m_r$ is a positive integer. Upon substituting $t_p=p^s$,
\eqref{dist:eq:dist-normal-coefficient} becomes the entire exponential polynomial
\begin{align}\label{conductor:eq:Aanalytic}
A_{d,f,\chi}(s)
&=\left(\frac df\right)^s
\prod_{\substack{p\mid d\\p\nmid f}}(1-\chi(p)p^{-s})\notag\\
&=\sum_{r\mid R(d,f)}\mu(r)\chi(r)m_r^s.
\end{align}
In particular,
\begin{equation}\label{conductor:eq:Aderivative}
A_{d,f,\chi}^{(j)}(s)
=\sum_{r\mid R(d,f)}\mu(r)\chi(r)m_r^s\log^j m_r.
\end{equation}

\begin{theorem}[All-index, all-parameter-order descent]\label{conductor:thm:stieltjes}
For $n\ge0$, $k\ge1$ and $f\mid d$,
\begin{align}\label{conductor:eq:stieldescent}
G_{d,\chi}^{(n,k)}
={}&\sum_{r\mid R(d,f)}\mu(r)\chi(r)m_r^{k+1}\notag\\
&\quad\cdot\sum_{j=0}^{n}(-1)^j\binom nj\log^j m_r\,
G_{f,\chi}^{(n-j,k)}.
\end{align}
No Euler-factor denominator occurs. Equivalently,
\begin{equation}\label{conductor:eq:stielA}
G_{d,\chi}^{(n,k)}=
\sum_{j=0}^{n}(-1)^j\binom nj A_{d,f,\chi}^{(j)}(k+1)
G_{f,\chi}^{(n-j,k)}.
\end{equation}
\end{theorem}
\begin{proof}
For $\Real s>1$, sorting a Dirichlet series by residue classes gives
\[
\sum_{a\in\U_d}\chi(a)\zeta(s,a/d)=d^sL(s,\chi_d),
\]
where $\chi_d$ is the induced character modulo $d$. Removing the extra
Euler factors from its primitive $L$-series yields
\[
L(s,\chi_d)=L(s,\chi)
\prod_{\substack{p\mid d\\p\nmid f}}(1-\chi(p)p^{-s}).
\]
Thus the character sum at $d$ is $A_{d,f,\chi}(s)$ times the character
sum at $f$. This is also the analytic realization of
Theorem~\ref{dist:thm:dist-normal-form}.

Insert $s=k+1+\eps$, multiply by the common factor
$(-1)^k(1+\eps)_k$, and use \eqref{tower:eq:masterN}. Comparison of
$(-1)^n\eps^n/n!$ gives \eqref{conductor:eq:stielA}; the sign $(-1)^j$ comes
from the Stieltjes Laurent convention. Finally apply
\eqref{conductor:eq:Aderivative}.
\end{proof}

\subsection{Concrete second-index identities at level twelve}

The multipliers are
\[
A_{12,3,\chi_3}(s)=4^s+2^s,\qquad
A_{12,4,\chi_4}(s)=3^s+1.
\]
Therefore, at $n=2$, $k=1$,
\begin{align}\label{conductor:eq:examplegamma}
G_{12,\chi_3}^{(2,1)}
&=20G_{3,\chi_3}^{(2,1)}
-72\log2\,G_{3,\chi_3}^{(1,1)}
+68\log^22\,G_{3,\chi_3}^{(0,1)},\\
G_{12,\chi_4}^{(2,1)}
&=10G_{4,\chi_4}^{(2,1)}
-18\log3\,G_{4,\chi_4}^{(1,1)}
+9\log^23\,G_{4,\chi_4}^{(0,1)}.
\end{align}
For instance,
\[
G_{12,\chi_3}^{(2,1)}=
\gamma_2'(1/12)-\gamma_2'(5/12)+\gamma_2'(7/12)-\gamma_2'(11/12),
\]
and $G_{3,\chi_3}^{(2,1)}=\gamma_2'(1/3)-\gamma_2'(2/3)$.
These are exact identities between explicitly specified signed sums; they
require no integer-relation search. The two identities were also checked
numerically using independent finite Hurwitz character sums.

\subsection{The undifferentiated layer and its pole correction}

For $k=0$, a principal character sum still contains the Hurwitz pole.
The same formula without a correction would be wrong.

\begin{proposition}[The $k=0$ correction]\label{conductor:prop:kzero}
For nonprincipal primitive $\chi$ the formula
\eqref{conductor:eq:stielA} remains valid with $k=0$. For $f=1$, put
$A_d(s)=A_{d,1,1}(s)$. Then
\begin{equation}\label{conductor:eq:kzero}
G_{d,1}^{(n,0)}=
\sum_{j=0}^{n}(-1)^j\binom nj A_d^{(j)}(1)\gamma_{n-j}(1)
+\frac{(-1)^n}{n+1}A_d^{(n+1)}(1).
\end{equation}
\end{proposition}
\begin{proof}
For a nonprincipal character, the sum of its values over units is zero,
so the $1/\eps$ residues cancel. Coefficient comparison is unchanged.
For the principal conductor, multiply
\[
A_d(1+\eps)\left(\frac1\eps+
\sum_{n\ge0}\frac{(-1)^n\gamma_n(1)}{n!}\eps^n\right).
\]
The pole term contributes $A_d^{(n+1)}(1)/(n+1)!$ to the coefficient of
$\eps^n$. Multiplying by $(-1)^n n!$ gives the correction in
\eqref{conductor:eq:kzero}. The residue itself is
$A_d(1)=d\prod_{p\mid d}(1-p^{-1})=\varphi(d)$, as required.
\end{proof}

For the non-character multiplication formula this specializes to
\begin{equation}\label{conductor:eq:classicalstieldist}
\sum_{j=0}^{d-1}\gamma_n\!\left(\frac{a+j}{d}\right)
=d\sum_{r=0}^{n}\binom nr(-\log d)^r\gamma_{n-r}(a)
+\frac{(-1)^n d\log^{n+1}d}{n+1}.
\end{equation}
The inhomogeneous term is the residue contribution, not an optional
normalization.


\section{Cyclotomic polylogarithm traces}\label{jets:sec:traces}

The same finite module applies to a second evaluation. Put
$\zeta_m=\exp(2\pi\iu/m)$ and evaluate
$e_x\mapsto\Li_s(\ee^{2\pi\iu x})$.
For $\Real s>1$, grouping the defining series by multiples of $p$ gives
\[
 \sum_{py=x}\Li_s(\ee^{2\pi\iu y})
 =p^{1-s}\Li_s(\ee^{2\pi\iu x}).
\]
Thus the weights are now $A_p=p^{1-s}$, and analytic continuation
extends the identity wherever the individual expressions are defined.

For a primitive character $\chi$ of conductor $f\mid q$, define
\begin{equation}\label{jets:eq:trace-definition}
 \cP_{q,\chi}(s)=
 \sum_{u\in U_q}\chi(u)\Li_s(\zeta_q^u),\qquad
 \tau(\chi)=\sum_{a\in U_f}\chi(a)\zeta_f^a.
\end{equation}
At conductor one take $\tau(\mathbf1)=1$ and $L(s,\mathbf1)=\zeta(s)$.

\begin{theorem}[Conductor-lifted trace identity]\label{jets:thm:trace}
As a meromorphic identity, and with removable singularities interpreted
by continuation,
\begin{equation}\label{jets:eq:trace-formula}
 \boxed{\displaystyle
 \cP_{q,\chi}(s)=
 (q/f)^{1-s}
 \prod_{\substack{p\mid q\\p\nmid f}}
 \bigl(1-\chi(p)p^{s-1}\bigr)
 \tau(\chi)L(s,\overline\chi).}
\end{equation}
For every $q>1$, the left side is an entire function of the spectral
variable $s$.
\end{theorem}
\begin{proof}
At primitive conductor $f>1$, absolute convergence and the primitive
Gauss-sum identity give
\[
 \sum_{a\in U_f}\chi(a)\Li_s(\zeta_f^a)
 =\sum_{n\ge1}\frac{\sum_a\chi(a)\zeta_f^{an}}{n^s}
 =\tau(\chi)L(s,\overline\chi).
\]
The Gauss-sum convention, including the complex conjugate on $\chi$,
is the one in~\cite{jets:dlmf-dirichlet}. Applying
\eqref{jets:eq:conductor-normal} with $A_p=p^{1-s}$ multiplies this
primitive trace by $E_{q,\chi}$. Extracting $p^{1-s}$ from each
first-appearance factor gives precisely the factor in
\eqref{jets:eq:trace-formula}. The conductor-one case follows directly
from $\Li_s(1)=\zeta(s)$ and the same polynomial identity.

For completeness, when $z=\zeta_q^u\ne1$, regrouping by residues gives
\begin{equation}\label{jets:eq:polylog-Hurwitz-DFT}
 \Li_s(z)=q^{-s}\sum_{a=1}^q z^a\zeta(s,a/q).
\end{equation}
Hurwitz zeta has only its simple pole at $s=1$, whose coefficient in
this finite sum is $\sum_a z^a=0$. Hence $\Li_s(z)$ is entire in $s$.
All arguments in the primitive level-$q$ trace differ from one.
\end{proof}

\begin{remark}[The character phase is consequential]
For the order-four character modulo five with $\chi(2)=\iu$, the
level-ten lifting factor is $2^{1-s}-\iu$, not
$2^{1-s}+\iu$. A complex-character numerical check is included to
catch an accidental interchange of $\chi$ and $\overline\chi$.
Real quadratic examples alone would not detect that error.
\end{remark}

\subsection{Exact nonprincipal vanishing orders}
For a nonprincipal primitive $\chi$, set
\[
 \rho_\chi=\#\{p\mid q:p\nmid f,\ \chi(p)=1\}.
\]
These are the same conductor statistics that controlled the
primitive-grid jet defects at Hurwitz order zero.

\begin{theorem}[Character trace zeros at order one]\label{jets:thm:trace-zero}
The entire function $\cP_{q,\chi}(s)$ has a zero of exact order
$\rho_\chi$ at $s=1$, where order zero means a nonzero value.
All lower derivatives vanish, and
\begin{align}\label{jets:eq:trace-leading}
 \cP_{q,\chi}^{(\rho_\chi)}(1)
 ={}&(-1)^{\rho_\chi}\rho_\chi!\,
 \tau(\chi)L(1,\overline\chi)\notag\\
 &\times\prod_{\substack{p\mid q,\ p\nmid f\\\chi(p)=1}}\log p
 \prod_{\substack{p\mid q,\ p\nmid f\\\chi(p)\ne1}}(1-\chi(p)).
\end{align}
\end{theorem}
\begin{proof}
Each factor with $\chi(p)=1$ in \eqref{jets:eq:trace-formula} is
$1-p^{s-1}=-(s-1)\log p+O((s-1)^2)$.
All other displayed factors are nonzero at $s=1$.
A primitive Gauss sum is nonzero, and Dirichlet's nonvanishing theorem
states $L(1,\overline\chi)\ne0$ for nonprincipal characters
\cite[Eq.~25.15.9]{jets:dlmf-dirichlet}. Thus there is no additional zero.
Taking the first nonzero Taylor coefficient proves the formula.
\end{proof}

Increasing the exponents of primes already present in $q$ changes
higher Taylor coefficients but not this first nonzero derivative:
$(q/f)^{1-s}$ equals one at $s=1$. This invariance provides an
additional family of identities between traces at different levels.

\subsection{The principal trace and its pole cancellation}
Write $\cP_q=\cP_{q,\mathbf1}$ and define
\[
 r=\omega(q),\quad P=\prod_{p\mid q}\log p,\qquad
 B=\log q-\tfrac12\log\rad(q).
\]
The conductor-one factor contains a pole, so the vanishing order is
one less than the number of primes, not equal to that number.

\begin{theorem}[Principal trace jet generator]\label{jets:thm:principal}
In a neighborhood of $t=0$,
\begin{equation}\label{jets:eq:principal-generator}
 \boxed{\displaystyle
 \cP_q(1+t)=(-1)^rP\,t^{r-1}\,
 \ee^{-Bt}\,[t\zeta(1+t)]
 \prod_{p\mid q}\frac{\sinh(t\log p/2)}{t\log p/2}.}
\end{equation}
In particular its zero at $s=1$ has exact order $r-1$, and
\begin{align}
 \cP_q^{(j)}(1)&=0 &&(0\le j<r-1),\label{jets:eq:principal-vanish}\\
 \cP_q^{(r-1)}(1)&=(-1)^r(r-1)!P,\label{jets:eq:principal-first}\\
 \cP_q^{(r)}(1)&=(-1)^r r!P(\gamma-B),\label{jets:eq:principal-second}\\
 \cP_q^{(r+1)}(1)&=(-1)^r(r+1)!P
 \left[-\gamma_1-B\gamma+\frac{B^2}{2}
 +\frac1{24}\sum_{p\mid q}(\log p)^2\right].
 \label{jets:eq:principal-third}
\end{align}
Here $\gamma_1$ is the ordinary first Stieltjes constant in
\eqref{jets:eq:stieltjes-laurent}.
\end{theorem}
\begin{proof}
The principal case of \eqref{jets:eq:trace-formula} is
\[
 \cP_q(1+t)=q^{-t}\prod_{p\mid q}(1-p^t)\zeta(1+t).
\]
Use $1-\ee^x=-x\ee^{x/2}\sinh(x/2)/(x/2)$ in every factor.
This proves \eqref{jets:eq:principal-generator}.
The factors in square brackets and in the hyperbolic-sine product are
analytic and equal to one at zero. Expanding them gives
\[
 t\zeta(1+t)=1+\gamma t-\gamma_1t^2+O(t^3),\qquad
 \prod_p\frac{\sinh(t\log p/2)}{t\log p/2}
 =1+\frac{t^2}{24}\sum_p(\log p)^2+O(t^4).
\]
Multiplication by $\ee^{-Bt}$ gives all the displayed derivatives.
\end{proof}

\subsection{Concrete exact identities}
The following examples are proved consequences of the preceding
theorems, not integer-relation conjectures. The derivatives are in the
\emph{order} of the polylogarithm, not its argument.

At level thirty,
\begin{align}
 \sum_{u\in U_{30}}\Li_1(\zeta_{30}^u)&=0,\label{jets:eq:q30-zero}\\
 \sum_{u\in U_{30}}\Liord{1}{\zeta_{30}^u}&=0,\label{jets:eq:q30-first}\\
 \sum_{u\in U_{30}}\Liord{2}{\zeta_{30}^u}
 &=\boxed{-2\log2\log3\log5}.\label{jets:eq:q30-second}
\end{align}
The third derivative of this trace is
$-6\log2\log3\log5\,[\gamma-\frac12\log30]$.
At level $210$, derivatives of orders zero, one, and two vanish, while
\begin{equation}\label{jets:eq:q210}
 \sum_{u\in U_{210}}\Liord{3}{\zeta_{210}^u}
 =6\log2\log3\log5\log7.
\end{equation}

For the quadratic character $\chi_3$ with $\chi_3(1)=1$ and
$\chi_3(2)=-1$, the primitive trace at order one is
$\tau(\chi_3)L(1,\chi_3)=\iu\pi/3$. Since
$\chi_3(7)=\chi_3(13)=1$, one obtains
\begin{align}
 \sum_{u\in U_{21}}\chi_3(u)\Liord{1}{\zeta_{21}^u}
 &=-\frac{\iu\pi}{3}\log7,\label{jets:eq:chi3-21}\\
 \sum_{u\in U_{273}}\chi_3(u)\Liord{2}{\zeta_{273}^u}
 &=\frac{2\iu\pi}{3}\log7\log13.\label{jets:eq:chi3-273}
\end{align}
In the first equation the order-zero trace vanishes; in the second,
both order-zero and first-order traces vanish.
Likewise the character $\chi_4$ gives
\begin{equation}\label{jets:eq:chi4-20}
 \sum_{u\in U_{20}}\chi_4(u)\Liord{1}{\zeta_{20}^u}
 =-\frac{\iu\pi}{2}\log5.
\end{equation}
These primitive trace constants follow directly from
$\Li_1(z)=-\log(1-z)$ with the values continued from inside the unit
disk, so the signs do not depend on an unspecified logarithm branch.

For the real quadratic character $\chi_5$, positive on residues
$1,4$ and negative on $2,3$, let $\phi=(1+\sqrt5)/2$.
Pairing conjugates in the primitive trace gives
\[
 \cP_{5,\chi_5}(1)
 =2\log\frac{\sin(2\pi/5)}{\sin(\pi/5)}=2\log\phi.
\]
As $\chi_5(11)=1$,
\begin{equation}\label{jets:eq:chi5-55}
 \sum_{u\in U_{55}}\chi_5(u)\Liord{1}{\zeta_{55}^u}
 =-2\log\phi\log11.
\end{equation}
This derivation needs no numerical integer-relation search or
unproved independence statement about logarithms.

\subsection{A stable independent formula for numerical checks}
Differentiation of \eqref{jets:eq:polylog-Hurwitz-DFT} at its cancelled pole
gives, for $z=\zeta_q^u\ne1$,
\begin{equation}\label{jets:eq:stable-polylog-jet}
 \Liord{k}{z}
 =\frac{(-1)^k}{q}\sum_{a=1}^q z^a
 \sum_{j=0}^k\binom kj(\log q)^{k-j}\gamma_j(a/q).
\end{equation}
There is no pole term, because its complete Fourier coefficient
vanishes identically. This supplies a finite Stieltjes-coordinate
evaluator for the trace identities, independently of their
Euler-factor right sides. In development, direct tiny-step numerical
differentiation of a polylogarithm routine near integral order gave
an unreliable trace derivative; the pole-cancelled formula avoids
that cancellation problem. This observation is an implementation
diagnostic, not a claimed theorem about all versions of a library.

\subsection{A level-260 specialization and convention reconciliation}
We use the unbarred weight $\chi(a)$ in $\mathcal P_{q,\chi}$.
An alternative convention weights the trace by $\overline\chi(a)$; it
conjugates the character factors and interchanges the primitive $L$-character.
All formulas here use the first convention. The following exact specialization
adds two primes with character value one.



Let $\chi$ be primitive of conductor $f>1$, and $f\mid q$.
Among the added primes, define
\[
\mathcal S=\{p\mid q:p\nmid f,\ \chi(p)=1\},\qquad r_\chi=|\mathcal S|.
\]
The remaining added primes form $\mathcal N$.

\begin{theorem}[Twisted leading jet]\label{conductor:thm:twisted}
The trace $\cP_{q,\chi}$ vanishes to order at least $r_\chi$ at one, and
\begin{equation}\label{conductor:eq:twistedlead}
\cP_{q,\chi}^{(r_\chi)}(1)=
(-1)^{r_\chi}r_\chi!\,
\cP_{f,\chi}(1)
\prod_{p\in\mathcal S}\log p
\prod_{p\in\mathcal N}(1-\chi(p)).
\end{equation}
Its order is exactly $r_\chi$ whenever $\cP_{f,\chi}(1)\ne0$.
More generally its exact order is $r_\chi+\ord_{s=1}\cP_{f,\chi}(s)$.
\end{theorem}
\begin{proof}
The multiplier is
\[
B_{q,f,\chi}(1+\eps)=
(q/f)^{-\eps}
\prod_{\substack{p\mid q\\p\nmid f}}(1-\chi(p)p^\eps).
\]
The factors indexed by $\mathcal S$ have simple zeros with leading
coefficient $-\log p$. All other factors are nonzero at zero and have
value $1-\chi(p)$. Multiply their leading terms by the Taylor series
of $\cP_{f,\chi}(1+\eps)$.
\end{proof}

For primitive nonprincipal characters, the classical nonvanishing
$L(1,\overline\chi)\ne0$, together with the nonzero primitive Gauss sum in
\eqref{jets:eq:trace-formula}, supplies the stated nonvanishing hypothesis
\cite{jets:dlmf-dirichlet}. The explicit example below does not need that theorem;
its primitive trace is computed directly from logarithms.

For $\chi_4(a)=(-1)^{(a-1)/2}$ on odd integers,
\[
\cP_{4,\chi_4}(1)=\Li_1(\iu)-\Li_1(-\iu)=\frac{\iu\pi}{2}.
\]
Take $q=260=4\cdot5\cdot13$. Both added primes have $\chi_4(p)=1$.
Consequently
\begin{align}\label{conductor:eq:260}
\sum_{a\in\U_{260}}\chi_4(a)\Li_1(\zeta_{260}^a)&=0,\\
\sum_{a\in\U_{260}}\chi_4(a)\Liord{1}{\zeta_{260}^a}&=0,\\
\sum_{a\in\U_{260}}\chi_4(a)\Liord{2}{\zeta_{260}^a}
&=\iu\pi\log5\log13.
\end{align}
This is a proved identity, not an integer-relation conjecture. More generally,
for distinct primes $p_1,\ldots,p_r\equiv1\pmod4$ and
$q=4p_1\cdots p_r$,
\begin{equation}\label{conductor:eq:chi4family}
\sum_{a\in\U_q}\chi_4(a)\Liord{r}{\zeta_q^a}
=(-1)^r r!\frac{\iu\pi}{2}\prod_{j=1}^r\log p_j,
\end{equation}
with all lower order traces vanishing. The same leading identity persists
if the added prime factors occur with higher exponents.

